%% kapa: exam=1 round=21 year=2010 period=1 form=A subject=economics subject_ko=경제원론 seq=33 number=73
%% source: Q-Net 감정평가사 기출문제 (cst00309) artlSeq=5024181 file=extracted/1008/5024181/2010년 제21회 감정평가사 1차시험 기출문제/1교시.hwp
%% conversion: kordoc 4.24.51 -> tools/gen1.py (auto); review=proofread-sonnet; flags=none
\begin{question}{73}
GDP(gross domestic product)에 관한 설명으로 옳지 않은 것은?
\begin{choices}
  \item GDP를 측정할 때 중간재의 가치는 제외하고 최종 상품과 최종 서비스의 가치만을 더한다.
  \item 실질 GDP의 단기 변화는 외국에서 자국민이 생산한 금액을 반영한다.
  \item 주부의 가사노동으로 생산된 금액은 GDP 계산에서 제외된다.
  \item 상품의 품질향상은 GDP 계산에 제대로 반영되지 못한다.
  \item GDP는 지하경제를 제대로 반영하지 못한다.
\end{choices}
\end{question}
\answer{2}
